\section{Full Proofs: Validity of the Convention, Sums, Certificates and Loopy Games}\label{S-kingsproof}
% Section S12, first subsection: the four results marked S in Sections 4 and 5 of the main text,
% with every definition they need and complete hand proofs. Results are numbered A.1, A.2, ... by hand.
\subsection{The Four Main Results, Self-Contained}\label{S-core}

This subsection states and proves the four results of the main text that carry the label S: Theorem~\ref{thm:validity},
Corollary~\ref{cor:v4}, the certificate statement of Section~\ref{sec:kings-valid} and Theorem~\ref{thm:hss}. It can be
read without the rest of this section. All definitions are given here, and every proof is a hand proof that uses only
these definitions and earlier steps of this subsection. No step rests on a computation. The one computed statement we
quote (part (c) of Proposition~A.14) is marked as such and is used nowhere. Numbered items are A.1, A.2, and so on.
We assume the basic theory of short games: disjunctive sums, outcome classes, and Zermelo's theorem that a finite game
without draws has a winner~\cite{Siegel2013}.

\subsubsection*{A.I Board, pawns and passed pawns}

\textbf{Board.} The board has $n$ files and $H\ge 3$ ranks. (On three ranks every pawn is passed, so the results below are about the empty position only; on four ranks a double step
promotes and no e.p.\ capture is possible.) It is \emph{flat} (files $1,\dots,n$ in a row) or a
\emph{cylinder} (file $n$ is also next to file $1$). The \emph{distance} of two files is their distance in the row, or
on the cycle for the cylinder. White pawns move up, black pawns down. A square is \emph{ahead} of a white pawn if its rank is higher than the pawn's, and ahead of a black pawn if
its rank is lower. No pawn stands on its own first rank; since pawns only move forward, this holds in all play. In this subsection a position holds pawns and at most two
kings, and no other piece until a promotion.

\textbf{Pawn moves.} A pawn steps one square forward to an empty square; from its own second rank it may instead
\emph{double step} two squares forward over an empty square to an empty square; it captures an enemy pawn diagonally
forward on a file at distance $1$ (across the seam on the cylinder). If the last move was a double step of a pawn $q$ from
$a$ over $b$ to $c$, an enemy pawn on the square beside $c$ (same rank, file at distance $1$) may capture \emph{en passant}
(e.p.): it moves to $b$ and $q$ is removed. This is allowed on the next move only. The position then carries a
\emph{right}; the right is \emph{usable} if such an enemy pawn exists. A pawn that reaches its last rank promotes; in the pawn games $W$ and $S$
below it stays there as a pawn, as in our engine. Kings are
never captured.

\textbf{A.1 Lemma (files).} \emph{A pawn changes file only by a capture, and it then lands on a file that held an enemy
pawn just before the move. Hence, in play in which kings are the only pieces other than pawns (all play in $W$, $S$, the
move bound and the threats below, where a promoted pawn stays a pawn, and chess before the first promotion), no pawn ever
enters a file that holds no pawn.}

\begin{proof}
A step or double step keeps the file. A capture lands on the square of the captured pawn; an e.p.\ capture lands on $b$,
on the file of the captured pawn $q$. Kings are never captured, so in such play the only capturable pieces are pawns.
\end{proof}

\textbf{Components.} A \emph{component} of a pawn position is a maximal set of consecutive files that all hold pawns (an
arc on the cylinder), with the pawns on them; on a cylinder with no empty file the whole board is one component, a
\emph{ring}. The files next to a component are empty. In play with pawns only, a move of a pawn in a component $D$ uses only squares on the files
of $D$ (Lemma~A.1), so it changes no other component, and whether it is legal depends only on $D$ and on the right.

\textbf{Passed pawns.} A pawn passes the \emph{textbook test} if no enemy pawn stands ahead of it on its own file or on a
file at distance~$1$. It passes the \emph{e.p.-aware test} if it passes the textbook test and it is not the case that it
has just double stepped with an enemy pawn beside it (Definition~\ref{def:epaware}). A \emph{setting} fixes the test
(textbook, e.p.-aware, or e.p.\ off: no e.p.\ captures and the textbook test) and a \emph{race rule}, which names a winner
when both sides have a passed pawn. We only use that the race rule is \emph{local}: its verdict depends only on the passed
pawns, the pawns on their files and on the files at distance~$1$, and on which side made the last move. The count, mover and
opponent rules of Section~\ref{sec:rules} are local. Under every test, whether a pawn of a component is passed depends only
on that component and on the right, since the files at distance $1$ from its file belong to the component or are empty.

\textbf{A.2 Lemma.} \emph{(a) A pawn on its last rank is passed under every test. (b) Let $\Phi_t$ be the sum, over the
pawns of side $t$, of the number of single steps the pawn needs to reach its last rank. Every
pawn move of $t$ lowers $\Phi_t$ by at least $1$, and no move raises $\Phi_t$. (c) Let $H\ge5$. If a position has no
passed pawn and its right, if any, was created by the last move of the side not to move, the side to move has no
promoting move.}

\begin{proof}
(a) No square lies ahead of it, so it passes the textbook test. It fails the e.p.\ clause only if it has just double
stepped with an enemy pawn beside it; that pawn would stand on its own first rank, where no pawn stands. (b) Each pawn move
advances the pawn by one or two ranks; a capture removes a pawn of the other side. (c) Since $H\ge5$, a double step
does not reach the last rank, so a promoting pawn stands on the rank before its last. Only the last rank lies ahead of
it, and it holds no enemy pawn, so the pawn passes the textbook test. The e.p.\ clause concerns only the pawn whose double
step created the right, and by hypothesis that pawn belongs to the side not to move, so the pawn is passed.
\end{proof}

\noindent Both restrictions in (c) are needed. On four ranks, in \pos{W a2 / B b3} no pawn is passed and a2--a4
promotes. On five ranks with the e.p.-aware test, after a2--a4 in \pos{W a2 b2 / B b4} no pawn is passed, yet a4--a5
promotes; White can play it only by moving twice in a row, which happens in a sum. Part (a) is what makes $W$ and
$S$ well defined in all these cases: a promoting move makes a passed pawn.

\subsubsection*{A.II The whole-board game, the sum model and the hypotheses of Theorem~\ref{thm:hss}}

We write White $=$ Left and Black $=$ Right, and $\bar\sigma$ for the opponent of a side $\sigma$. A positive integer is
the game $k=\game{k-1}{}$, with $0=\game{}{}$: Black has no move in it, and White has exactly $k$ moves in it before it is
$0$. We write $\sigma\BIG$ for $+\BIG$ if $\sigma$ is White and $-\BIG$ if $\sigma$ is Black.

\textbf{Whole-board game $W(X)$.} Fix a setting. Let $X$ be a pawn position, possibly with a right, that has no passed
pawn. A position of $W(X)$ is a pawn position with its right and the side to move. A move is a legal pawn move. If after a
move some side has a passed pawn, the game ends: if only $\sigma$ has one, $\sigma$ wins; if both have one, the race rule
names the winner. Otherwise play goes on, and a player with no move loses. By Lemma~A.2(b), $\Phi_W+\Phi_B$ falls at every
move, so $W(X)$ is finite and has no draws; by Lemma~A.2(a) a promoting move ends it. We write $\mathrm{Reach}(X)$ for the
positions of $W(X)$ reached from $X$ by play, with either side to move first, before the game ends ($X$ included). The
game $W(X)+\nim{h}$ is played in the same way, with one more move available: a move in the heap. A passed pawn ends the whole
game with the same winner, and a player with no move in $W(X)$ and an empty heap loses. (Replacing the end by an integer
$\pm\BIG$ with $\BIG>0$ and continuing in the heap gives the same outcomes, since $\BIG+\nim{h}$ is won by White whoever
starts: White empties the heap or moves in $\BIG$, and Black has no move in a nonnegative integer.)

\textbf{Sum model $S(X)$.} For a component $D$ (with its right, if the right concerns its pawns) let
\[
S(D)=\game{S(D^L_1),\,S(D^L_2),\dots}{S(D^R_1),\dots},
\]
where $D^L_i$ runs over the positions on the files of $D$ after a legal white move in $D$ (with the right that move
creates), $D^R_j$ likewise for black, and for such a position $D'$ we set $S(D')=\sigma\BIG$ if some side has a passed pawn
in $D'$ and $\sigma$ is the winner as in $W$, and otherwise $S(D')=\sum_E S(E)$ over the components $E$ of $D'$. Then
$S(X)=\sum_D S(D)$ over the components of $X$. The recursion ends by Lemma~A.2(b), and by Lemma~A.2(a) a promoting move
freezes its component. This is the main text's $S(X)$. Note the order: after a move the whole moved component is tested first, and
frozen at $\pm\BIG$ if a side has a passed pawn in it; only otherwise is it re-split. (Splitting first would be a different model: in \pos{W b2 d6 / B c7} on the ordinary board, c7$\times$d6 leaves two passed pawns on separated files, a race that White wins under the count rule, whereas $+\BIG-\BIG=0$.) In a sum the right of a component
stays with that component while the players move elsewhere; in $W(X)$ a right lapses after one move. This difference is
what condition (E) controls.

\textbf{Move bound.} $N(X)$ is the largest number of moves of one side in a finite sequence of legal pawn moves from $X$,
made by either side in any order, where e.p.\ is allowed only directly after the double step it answers and play goes on
past passed pawns. By Lemma~A.2(b), $N(X)\le\max(\Phi_W(X),\Phi_B(X))$, and a sequence from a position of
$\mathrm{Reach}(X)$, prefixed by the play that reached it, is a sequence from $X$.

\textbf{Threats, HSS and (E).} A component $D$ with no passed pawn \emph{carries a threat for $t$} if there is a finite
sequence of legal pawn moves on the files of $D$, starting from $D$ with its right, made by either side in any order, such
that no position before the last has a passed pawn and the last position gives $t$ a passed pawn (a race gives both sides
one). The right of each position in the sequence is the one created by the move before it.
$X$ is \emph{hereditarily sum-safe} (HSS) if no $Y\in\mathrm{Reach}(X)$ has two components that carry threats for opposite
sides. $X$ satisfies \emph{(E)} if no $Y\in\mathrm{Reach}(X)$ with at least two components carries a usable right.

\textbf{A.3 Remark (the precise form is implied by the main text's).} Definition~\ref{def:hss} and Theorem~\ref{thm:hss}
quantify over positions reachable by moves of either side \emph{in any order}; we use only $\mathrm{Reach}(X)$, which is
smaller, so our HSS and (E) are implied by the main text's. The main text's threat is the same notion in words: play on the component's own files, with no passed
pawn earlier (a race gives both sides one). Dropping the clause about earlier passed pawns would give more threats and a stronger
hypothesis; our reading has fewer threats, so the theorem below covers both readings. ``$\BIG$ exceeds the number of moves either side can make'' is $\BIG>N(X)$. The main
text leaves implicit that the race rule is local, and that a pawn on its last rank stays a pawn and so counts as passed, as in our engine.

\textbf{A.4 Lemma (what (E) gives).} \emph{Let $X$ satisfy (E), let $Y\in\mathrm{Reach}(X)$ have at least two components,
and let $\mu$ be a move in a component $C$ of $Y$, giving $Y'$. For every other component $D$ of $Y$: (a) the games
$S(D)$ with the right of $Y$ (if it concerns $D$) and without it have the same options; (b) no pawn of $D$ is passed in
$Y'$.}

\begin{proof}
(a) By (E) the right of $Y$ is not usable, so it allows no move, and every other move of $D$ leads to the same position
with the same new right. (b) Pawns of $D$ are not passed in $Y$, since the game did not end there. A pawn of $D$ that
passes the textbook test in $Y$ would have failed only the e.p.\ clause, so $Y$ would carry a usable right, against (E). So
no pawn of $D$ passes the textbook test in $Y$, and $D$ is unchanged in $Y'$.
\end{proof}

\textbf{A.5 Lemma (threats are inherited).} \emph{Let $E$ be a component with no passed pawn that carries no threat for
$\bar\sigma$. Let $\bar\sigma$ make a move $\mu$ in $E$, giving $E'$ on the files of $E$. Then $S(E')=\sigma\BIG$, or $E'$ has no passed pawn and each of its components carries no threat for
$\bar\sigma$.}

\begin{proof}
If $\bar\sigma$ has a passed pawn in $E'$ (alone or in a race), the one-move sequence $\mu$ shows that $E$ carries a
threat for $\bar\sigma$. So either only $\sigma$ has one, and $S(E')=\sigma\BIG$, or $E'$ has none. In the second case
suppose a component $E_i$ of $E'$ carries a threat for $\bar\sigma$, by a sequence $s$. We show that $\mu$ followed by $s$
witnesses a threat of $E$ for $\bar\sigma$. The moves of $s$ use only the files of $E_i$, so they are legal after $\mu$. If
$\mu$ is a double step, it changes no file's occupancy, so $E'$ has the files of $E$ and $E_i=E'$, with the right of $\mu$,
which is the right the sequence $s$ starts from. If $\mu$ is not a double step, $E'$ carries no right, and every other
component of $E'$ stays as it is in $E'$, with no passed pawn and no right, during $s$. In both cases every position before
the last has no passed pawn ($E$, $E'$, and those of $s$), and the last one gives $\bar\sigma$ a passed pawn. This
contradicts the hypothesis on $E$.
\end{proof}

\textbf{A.6 Lemma (frontier).} \emph{Let $X$ satisfy (E) and $\BIG>N(X)$. Let $Y\in\mathrm{Reach}(X)$, and let a move
$\mu$ in a component $C$ of $Y$ end $W$ with winner $\sigma$. Suppose no component $D\ne C$ of $Y$ carries a threat for
$\bar\sigma$. Then $\sigma$ wins $G=\sigma\BIG+\sum_{D\ne C}S(D)$ whoever moves first.}

\begin{proof}
Take $\sigma$ White; the other case is the mirror image. We play on the game trees as displayed, without replacing a pawn
component by a simpler game of the same value. White's strategy is to move, at each turn, in the original summand $+\BIG$, one move at a time.
Black has no move in a nonnegative integer, so every Black move is a pawn move in a summand $S(E)$, where $E$ is
a component of $Y$ other than $C$ or arises from one. By Lemma~A.5 and induction on the number of Black moves, each such $E$
has no passed pawn and carries no threat for Black, and every component that Black's move freezes becomes $+\BIG$. Black never captures
e.p.: White makes no pawn move after $\mu$, and the only right White may have created earlier is the right of $Y$, which
is not usable by (E), since $Y$ has the two components $C$ and $D$ (if $Y$ has no component $D\ne C$, Black has no pawn move
at all). So Black's moves, in order, follow the play from $X$ to $Y$ and $\mu$ as a sequence of legal pawn moves, and Black
makes at most $N(X)<\BIG$ of them. White needs at most one move per Black move, plus one if White starts, so at most
$N(X)+1\le\BIG$ moves, and the summand $\BIG$ alone supplies $\BIG$ of them. So White always has a move. If Black makes $k$ moves in all, White needs $k$ moves when Black starts and $k+1$ when White starts, and $k+1\le N(X)+1\le\BIG$. The play is finite
(the integers shrink and $\Phi_B$ falls), so it ends with a player who cannot move, and this is Black.
\end{proof}

\textbf{A.7 Theorem (sum safety; Theorem~\ref{thm:hss}).} \emph{Main text: ``Suppose that $X$ has no passed pawn and is HSS, that no position
reachable from $X$ with two or more components has a usable en passant right (condition (E)), and that $\BIG$ exceeds the
number of moves one side can make in any sequence of pawn moves from $X$. Then $S(X)$ and $W(X)$ have the same outcome.'' Precise form: let $X$ be a pawn
position with no passed pawn, fix a setting with a local race rule, and suppose that $X$ is HSS, satisfies (E), and
$\BIG>N(X)$. Then for each side to move first, $S(X)$ and $W(X)$ have the same winner.}

\begin{proof}
For $Y\in\mathrm{Reach}(X)$ let $S_Y=\sum_D S(D)$ over the components of $Y$, the right of $Y$ going with the component
of the pawn that made the double step. Then $S_X=S(X)$. A position of $\mathrm{Reach}(X)$ other than $X$ comes with the side $t$ to move in the play that reached it; at $X$ either side may be to move. We show, by induction on $\Phi_W+\Phi_B$, that $W$ at $Y$ and
$S_Y$ have the same winner with $t$ to move. (The other side to move at $Y$ is not needed and need not satisfy HSS or (E).)

The moves of $S_Y$ are the moves in the components of $Y$. These are exactly the legal moves of $W$ at $Y$: pushes and
captures use the files of one component (Lemma~A.1), and an e.p.\ capture uses the right of $Y$, which lies in the
component of the pawn that double stepped, next to the capturer. Fix a move $\mu$ in $C$, giving $Y'$ in $W$ and $C'$ on
the files of $C$. The option of $S_Y$ is $\sum_{D\ne C}S(D)+S(C')$.

\emph{$Y'$ has no passed pawn.} Then $C'$ has none, and $S(C')$ is the sum over the components of $C'$. The components of
$Y'$ are those of $C'$ and the $D\ne C$. If $Y$ has two or more components, Lemma~A.4(a) lets us drop the right of $Y$
from the summands $S(D)$ without changing them; no summand then keeps an older right. (If $Y$ has one component, $\mu$ is made in it and the right of $Y$ lapses in $S$ as in $W$.) So the option is $S_{Y'}$, and the induction hypothesis applies.

\emph{$Y'$ ends $W$ with winner $\sigma$.} In $W$, $\sigma$ wins at once. By Lemma~A.4(b) (or because $C$ is the only
component) every passed pawn of $Y'$ lies in $C'$, so $C'$ has a passed pawn, and $S(C')=\sigma\BIG$ because the race rule
is local and the mover is the same. The move $\mu$ alone shows that $C$ carries a threat for $\sigma$ (and for both sides
in a race). By HSS no other component of $Y$ carries a threat for $\bar\sigma$, and Lemma~A.6 shows that $\sigma$ wins the
option whoever moves first.

So every option has the same status in both games: won by the same player when the opponent moves next. A player with no
option loses in both. Hence the winners agree at $Y$, and at $Y=X$ the claim is the theorem.
\end{proof}

\textbf{Remark ((E) cannot be dropped).} On the flat board with $5$ files and $7$ ranks and e.p.\ on, let $X$ be
\pos{W a4 b3 e2 / B a5 b4 e4} with the right created by a2--a4. No pawn is passed, and the components are the files a, b
and the file e. The two pawns on e face each other and never pass, so that component carries no threat in any reachable
position and $X$ is HSS. The right is usable (b4 takes a4 e.p.\ on a3), so (E) fails. In $W(X)$ White, moving first, plays
e2--e3; the right lapses and Black has no move. Black, moving first, captures e.p.\ and has the only passed pawn. So the
first player wins $W(X)$. In $S(X)$ the right stays with the component on a, b, where White has no move and Black's
capture gives $-\BIG$; that component is $\{\ \mid-\BIG\}=-\BIG-1$, and the file e is $\nim{1}$. So Black wins $S(X)$
whoever starts. No race occurs and the argument uses neither passed test, so this holds in every setting with e.p.\ on
\lab{P}; the engine and a separately written minimax give the same two outcomes.

\textbf{A.8 Proposition (opposite threats cancel).} \emph{For every integer $\BIG\ge1$,
$\game{\BIG}{0}+\game{0}{-\BIG}=0$, although in the game where a $\pm\BIG$ ends play the first player wins.}

\begin{proof}
If White starts and moves to $\BIG+\game{0}{-\BIG}$, Black answers $-\BIG$, and White faces $0$; if White moves to
$\game{\BIG}{0}+0$, Black answers $0$. Either way White loses. Black starting loses in the same way. A second-player win
equals $0$. If instead reaching $\pm\BIG$ ends play, the first player moves to its own $\pm\BIG$ and wins.
\end{proof}

\textbf{A.9 Remark (the move bound).} A file contributes at most $1+\dots+(H-2)=(H-1)(H-2)/2$ to $\Phi_t$, when its squares on ranks $2$ to $H-1$ all hold pawns of $t$, so
$N(X)\le\max_t\Phi_t(X)\le n(H-1)(H-2)/2$, which is at most $168<256$ on boards with $n\le8$ files and $H\le8$ ranks. So $\BIG=256$ meets
$\BIG>N(X)$ there, as the main text says.

\subsubsection*{A.III King groups, the promotion game and the validity theorem}

\textbf{King group.} A \emph{king group} $Z$ is a set $F_Z$ of consecutive files (an arc on the cylinder) holding both
kings and some pawns, the \emph{Z-pawns}. The files at distance $1$ from $F_Z$ form the \emph{buffer} and hold no pawn.
The other pawns, the \emph{context}, stand on files at distance at least $2$ from $F_Z$. A \emph{Z-move} is a move of a
king or of a Z-pawn. $X$ is the pawn position of the context alone; we assume it has no passed pawn. In this part $H\ge5$, so that Lemma~A.2(c) applies.

\textbf{Three games.} \emph{Chess} $\mathrm{C}$ is played with the usual rules, any promotion, no fifty-move rule;
checkmate wins and stalemate and infinite play are draws. A side \emph{wins} a position if it has a strategy that ends every
play in a win for it in finitely many moves. The \emph{promotion game} $\mathrm{P}$ is chess in which the first promotion
ends the game and wins for the promoter. The \emph{frozen-king game} $\mathrm{F}$ is $\mathrm{P}$ in which, in addition,
a king move loses at once for its mover; our solver uses it to play the convention, and in it every king move is a loss by
definition.

\textbf{A.10 Lemma (buffer).} \emph{Before the first promotion, in any play from a position with a king group, no pawn
stands on a buffer file, Z-pawns stay on $F_Z$, context pawns stay on files at distance at least $2$ from $F_Z$, no context
pawn attacks a square of $F_Z$, and a Z-pawn captures only on $F_Z$.}

\begin{proof}
By Lemma~A.1 a pawn enters a file only if an enemy pawn stands on it. The buffer starts empty, so by induction on the moves
it stays empty. To leave $F_Z$, or to reach $F_Z$ from the context, a pawn would have to enter a buffer file. A pawn
attacks only the files at distance $1$ from its own, and a Z-pawn's capture lands on the file of the captured pawn (also in an e.p.\ capture), which is not a buffer file.
\end{proof}

\textbf{Decisive events.} Let $Q$ be a position made of a king group $Z$ and its context, with a side to move. A move is a
\emph{decisive event} if it is a Z-move, or a context move that leaves some side with a passed
pawn among the context pawns. Since the Z-pawns are at least two files away, this is also the passed-pawn test on the whole
board for the context pawns. $\mathrm{R}(Q)$ is the set of chess positions reached from $Q$ by legal moves none of which is
a decisive event ($Q$ included). Kings and Z-pawns do not move along such play, so a position $Y\in\mathrm{R}(Q)$ is fixed
by its context part $X_Y$, its right and its side to move.

\textbf{Conditions.} At every $Y\in\mathrm{R}(Q)$, with side $s$ to move:
\begin{itemize}
\item[(V1)] the legal chess moves of $s$ that are not Z-moves are exactly the moves of $s$ in $W$ at $X_Y$, with the same
resulting context and right; none of these moves gives check and nobody is in check; and $s$ has a legal chess move (a Z-move counts);
\item[(V2)] every context move at $Y$ that produces a passed pawn and is adjudicated for $\sigma$ in $W$ leads to a chess position that $\sigma$ wins;
\item[(V3)] every legal Z-move at $Y$ leads to a chess position that the opponent of its mover wins.
\end{itemize}
(V4) is: $X$ is HSS, satisfies (E), and $\BIG>N(X)$.

\textbf{A.11 Theorem (validity of the convention; Theorem~\ref{thm:validity}).} \emph{For the given side to move, the main text identifies the chess result with the outcome of $W(X)$ under (V1)--(V3). Precise form: let $Q$ be a
king group with its context, with no passed pawn in $X$, and let (V1)--(V3) hold at every position of $\mathrm{R}(Q)$. If
$\sigma$ wins $W(X)$ with the side to move of $Q$, then $\sigma$ wins $Q$ in chess. In particular $Q$ is not a draw.}

\begin{proof}
$W(X)$ is finite without draws, so some $\sigma$ wins it; fix a winning strategy. We give $\sigma$ a chess strategy and
show that every play under it is a $\sigma$-win. Along the play, as long as no decisive event has occurred, the position
$Y$ lies in $\mathrm{R}(Q)$ and $X_Y$ is a position of $W(X)$ won by $\sigma$ with the same side to move. This holds at $Q$.

\emph{$\sigma$ to move.} $\sigma$ plays the move of its strategy at $X_Y$. By (V1) it is a legal chess move with the same
result. If it produces a passed pawn, the winner is $\sigma$, and by (V2) the chess position reached is won by $\sigma$; $\sigma$ follows
a winning strategy from there. Otherwise the move is not decisive, the new position is in $\mathrm{R}(Q)$, and $W$ is won by
$\sigma$ there.

\emph{$\bar\sigma$ to move.} By (V1), $\bar\sigma$ has a legal move. A Z-move leads to a position won by $\sigma$, by (V3).
Any other legal move is, by (V1), a move of $W$ at $X_Y$, and all of these lead to positions won by $\sigma$ in $W$. If it
produces a passed pawn, the winner is $\sigma$, and (V2) applies. Otherwise the new position is in $\mathrm{R}(Q)$ and won by $\sigma$ in
$W$. (If $\bar\sigma$ has no move in $W$, its legal moves are all Z-moves.)

Every move before the first decisive event is a pawn move of $W$, so by Lemma~A.2(b) there are at most
$\Phi_W(X)+\Phi_B(X)$ of them; (V1) then still supplies a legal move, so a decisive event occurs by move $\Phi_W(X)+\Phi_B(X)+1$, and at the event (V2) or (V3) gives $\sigma$ a strategy that wins in finitely many further moves. Before the event, (V1) excludes stalemate; from the event onward the chosen winning continuation excludes it. No promotion occurs before the event: a Z-pawn promotion is a Z-move, and a context
promotion would make a passed pawn in the pawn-game successor by Lemma~A.2(a) and so would itself be an event. Moreover at each $Y\in\mathrm{R}(Q)$ the context
$X_Y$ has no passed pawn and its right, if any, was made by the side not to move, so by Lemma~A.2(c) no context move at $Y$
promotes; the moves compared in (V1) and (V2) are pawn moves that leave pawns.
\end{proof}

\textbf{A.12 Remark.} (a) The proof uses from (V1) only the move correspondence with matching successors and the existence of a legal move, not the
clauses on checks. A context move never gives check before a decisive event, by Lemma~A.10, since the kings stay on $F_Z$;
and a side in check from a Z-pawn has no legal context move, so the equality of move sets then fails unless $W$ has no move
either. (b) The theorem and the proof hold word for word
with chess replaced by $\mathrm{P}$ or by $\mathrm{F}$; in $\mathrm{F}$, (V3) holds for every king move by definition.
(c) $X$ must have no passed pawn; otherwise $W(X)$ is not defined. The main theorem explicitly assumes at least five ranks, as required for the use of Lemma~A.2(c).
(d) (V1) contains no condition (E): it compares chess with the whole-board game $W(X)$, whose e.p.\ rights lapse after one move as in chess. Condition (E) is needed only to pass from $W(X)$ to $S(X)$ (Theorem~A.7), hence in (V4).

\textbf{A.13 Corollary (Corollary~\ref{cor:v4}).} \emph{Main text: ``If also (V4) holds, that is, $X$ is hereditarily sum-safe, satisfies condition (E) of Theorem~\ref{thm:hss}, and $\BIG$ exceeds the number of moves either side can make, then for that side to move the chess result is the one given by the components' $\pm\BIG$ value.'' Precise form: for a side to move for which the hypotheses of Theorem~A.11 hold, assume also (V4) and a local race rule, as in Theorem~A.7. Then the chess winner of $Q$ with that side to move is the winner of $S(X)$ with that side moving first. (To infer the full outcome class one must verify the chess hypotheses for both starters. They can hold for one side only: on the ordinary board, with the king group \pos{W Kb6 c7 / B Ka8 a5} and no other pawn, they hold with Black to move, who must play a5--a4 and is mated by c8=Q, but not with White to move, who mates at once although the empty sum is lost for the side to move.)}

\begin{proof}
By Theorem~A.11 the chess winner is the winner of $W(X)$, and by Theorem~A.7 this is the winner of $S(X)$.
\end{proof}

\textbf{A.14 Proposition (Proposition~\ref{prop:sumcounter}).} \emph{On $\cyl{8}$ with e.p.\ off, the count race rule and $\BIG=256$, let
$X=$\,\pos{W b2 d4 g3 h5 / B a7 c7 f7 g4}. (a) $X$ is one component. (b) $X$ is not HSS. (c) (Checked by two separately written programs \lab{C2}, not
proved here and not used here.) $S(X)$ is a black win whoever starts and $W(X)$ a first-player win.}

\begin{proof}
(a) Only file e is empty. (b) White h5--h6 and Black a7$\times$h6 (across the seam) give
$Y=$\,\pos{W b2 d4 g3 / B c7 f7 g4 h6}. No pawn in $X$, after h5--h6, or in $Y$ is passed: g3 and g4 block each other, f7
has g3 ahead on a neighbouring file, c7 has b2 ahead, b2 and d4 have c7 ahead, a7 has b2 ahead, the white h-pawn has a7
ahead, and the black pawn on h6 has g3 ahead.
Files a and e of $Y$ are empty, so $Y$ has the components $\{f,g,h\}$ and $\{b,c,d\}$. The first carries a threat for
Black: f7--f6, f6--f5, f5--f4 keep g3 ahead of the f-pawn, and f4$\times$g3 leaves no white pawn in the component, so the new g3 pawn is passed. The second carries
one for White: d4--d5, d5--d6 keep c7 ahead of the d-pawn, and c7$\times$d6 empties file c, so b2 is passed (d6 is passed too, a race, which is a threat for both sides).
\end{proof}

\subsubsection*{A.IV The certificate}

\textbf{Data.} Let $Y$ be a position with a king group $Z$ and its context, no promotion yet, and side $s$ to move; write
$\bar s$ for the other side. Let $m$ be a legal Z-move of $s$ that does not promote (it may take $s$'s king off $F_Z$), and
$r$ a legal reply of $\bar s$ that is a Z-move, does not promote, and leaves $\bar s$'s king on $F_Z$. Let $Q$ be the
position after $m$ and $r$, with $s$ to move, and let $h\ge1$. The \emph{projection} $\pi(R)$ of a position $R$ deletes
every pawn off $F_Z$ and keeps the kings, the Z-pawns, the side to move, and a right whose captured pawn is on $F_Z$.

\textbf{Certificate of horizon $h$.} The move $m$ has one, with reply $r$, if (C1)--(C3) hold.
\begin{itemize}
\item[(C1)] \emph{Local promotion.} In the \emph{local game} from $\pi(Q)$, with $s$ to move, played by the rules of chess
except as follows, $\bar s$ has a strategy that wins against every play of $s$: $s$ may pass whenever it is not in check
(a pass changes only the side to move and clears the right); $\bar s$ may not pass and its king may only step onto $F_Z$;
the first promotion ends the game, a win for $\bar s$ if $\bar s$ promotes and a failure for $\bar s$ if $s$ does; a
checkmate of $s$ is a win for $\bar s$; $\bar s$ having no permitted move is a failure; and $\bar s$ must win with one of
its first $h$ moves, or it fails.
\item[(C2)] \emph{Counter-promotion.} Let $Q^\circ$ be $Q$ with $\bar s$'s king and all Z-pawns deleted. In $Q^\circ$
only $s$ moves: its king steps to squares that hold no pawn of $s$ and are not attacked by a pawn of $\bar s$ (it may
capture such pawns), and its pawns move by the pawn rules without e.p. Every sequence of such moves that ends with a
promotion has at least $h+1$ moves.
\item[(C3)] \emph{No stalemate.} There is a strategy for $\bar s$ from
$Q$ in $\mathrm{P}$ with the following property. Until the game ends or
$\bar s$ has made $h$ moves from $Q$, each move it prescribes is a legal
Z-move whose projection is winning in (C1) with the remaining number
of moves. In every such play $s$ is never stalemated before $\bar s$'s
$h$-th move. Call this a guided strategy satisfying (C3).
\end{itemize}

\textbf{A.15 Remark (the main text's form).} The main text says that $m$ has a certificate of horizon $h$ ``if, with
only $Z$ on the board, the punisher promotes or mates within $h$ moves of its reply, no defender pawn outside $Z$ is that
fast, and no stalemate comes first''. The quoted sentence summarizes (C1)--(C3). The clock counts the
punisher's moves after its reply. Condition (C1) uses the projected board,
a defender allowed to pass out of check, and a punisher whose king stays
on $F_Z$. Condition (C2) lets the defender move its king and context
pawns consecutively in $Q^\circ$, including king captures of blockers.
Condition (C3) requires a choice of guided real strategy that avoids
defender stalemate throughout the window. A clause that an unsafe queen is a failure, used in the long version with chess in view, only strengthens
(C1); it gives no chess win either (Remark~A.17).

\textbf{A.16 Theorem (certificate).} \emph{Main text: ``In the promotion game (the first promotion wins) such a move loses
for its mover.'' Precise form: if $m$ has a certificate of horizon $h$ with reply $r$, then after $m$ the side $\bar s$ wins
$\mathrm{P}$: it plays $r$ and wins by promotion or mate no later than its $h$-th
subsequent move. The reply itself may already mate; otherwise it follows
the guided strategy of (C3) for as long as play continues.}

\begin{proof}
Consider any play from $Q$ in which $\bar s$ follows the guided strategy $\rho$ and $s$ makes legal moves, up to $\bar s$'s
$h$-th move or the first promotion. Let $R$ denote its positions.

\emph{Invariants.} (I1) By Lemma~A.10 (from $Y$; the kings' squares play no role in it), the buffer has no pawn, Z-pawns
stay on $F_Z$, context pawns stay off $F_Z$ and the buffer, no context pawn attacks a square of $F_Z$, and Z-pawns capture
only on $F_Z$. (I2) $\bar s$ makes only Z-moves and its king stays on $F_Z$. (I3) $\bar s$'s context pawns never move and
are only removed. No $\bar s$ pawn of the context attacks $s$'s king: none does after $m$, since $m$ is legal; $r$ moves no
context pawn; later the set of squares these pawns attack only shrinks, and $s$'s king never steps onto an attacked square.
So $s$ is in check in $R$ exactly when a Z-pawn of $\bar s$ attacks its king, which is the same in $\pi(R)$; the kings
never attack each other. (I4) Rights: a right created by $\bar s$ comes from a Z-pawn and its captured pawn is on $F_Z$; a
right created by a context double step of $s$ could be used only by a $\bar s$ context pawn, which never moves. So the only
e.p.\ captures in the play are by Z-pawns of pawns on $F_Z$, and $\pi$ keeps those rights.

\emph{Step 1: moves of $s$ project to local moves or passes.} Let $s$ be to move in $R$. A king move of $s$ is legal in
$\pi(R)$: the target square holds the same piece or less, the kings are the same, and the pawns of $\bar s$ that attack it
in $\pi(R)$ also do in $R$; a capture of a context pawn becomes a plain move. A Z-pawn move of $s$ uses squares of $F_Z$,
which carry the same pieces in $R$ and $\pi(R)$, and the checks are the same by (I3); a Z-pawn promotion is a promotion in
the local game. A non-promoting context-pawn move leaves $\pi(R)$ unchanged except for the side to move and the right, as a pass does; it
is legal only if $s$ is not in check, since a context pawn can neither capture on $F_Z$ (I1) nor block a pawn check; so the
pass is allowed in $\pi(R)$. Hence, up to a possible first context promotion, each nonterminal real prefix projects to a legal prefix of the local game in which $\bar s$ plays only
winning moves, with the same count of moves of each side.

\emph{Step 2: permitted local moves of $\bar s$ are legal.} Let $\bar s$ be to move in $R$ and let $\nu$
be a permitted local move at $\pi(R)$. Its squares lie on $F_Z$, which holds the same pieces in $R$ and $\pi(R)$, and every right usable by this local move agrees in the two positions (I4). After $\nu$, $\bar s$'s king stands on $F_Z$, where no context pawn of $s$ attacks (I1); the other
attackers are the same as in $\pi(R)$. There are no sliding pieces, so there are no pins. So $\nu$ is legal in $R$. By (C1) at the root and induction using Step~1 and the previous guided choices, the live projected position is one that $\bar s$ wins in the local game with the moves that remain, so a winning move exists there, and the
guided strategy is defined at every turn of $\bar s$ before it wins; in particular $\bar s$ is never mated or stalemated.

\emph{Step 3: $s$ does not promote first.} In the window from $Q$ through $\bar s$'s $h$-th subsequent move, $s$ has at most $h$ turns, since $s$ moves first. A promotion of a Z-pawn of $s$ would be a promotion by $s$ in the
local game, a failure for $\bar s$, which cannot happen from a local position that $\bar s$ wins. Suppose a context pawn of $s$ promotes at
$s$'s $j$-th turn, $j\le h$. Delete from the play all moves of $\bar s$ and all Z-pawn moves of $s$. What remains is at most
$j$ moves of $s$, king moves and context-pawn moves, ending with the promotion. It is a legal sequence in $Q^\circ$: $s$'s
king and the context pawns stand on the same squares in both, since only $s$ moves them; a king move in the play goes to a
square not attacked by a context pawn of $\bar s$ and not holding a pawn of $s$; a context-pawn move uses context squares,
whose occupants are the same, since $\bar s$'s king and the Z-pawns are on $F_Z$; and no e.p.\ capture is involved (I4).
This contradicts (C2).

\emph{Step 4: $\bar s$ wins.} If $s$ is already mated at $Q$ or is
mated earlier on the real line, $\bar s$ has won. Otherwise (C3)
excludes real stalemate of $s$ at each of its turns before the deadline.
Step~2 supplies a legal winning move at each live punisher turn, and
Step~3 excludes defender promotion in the window. Thus while the real
line continues its projection is a local prefix with only winning
punisher moves at their remaining budgets. Such a prefix cannot continue
past the $h$-th punisher move without local promotion or mate. A local
punisher promotion is the same legal real promotion. At a local mate
$s$ is in check in the real game by (I3), and any real escape would
project to a legal local escape by Step~1. Hence the real game ends in
a punisher win by the stated deadline.
\end{proof}

\textbf{Remark (timing).} The reply $r$ itself may already mate while the projected position is still live: on the flat board take $F_Z=\{a,b\}$, buffer c, \pos{W Ka6 b6 b7 d6 / B Ka8 a7 e7} with Black to move, $m=\mathrm{Kb8}$ and $r=\mathrm{b6{\times}a7}$ mate. In $\pi(Q)$ Black has Kc7 and White promotes a8=Q, so $h=1$ works, yet the real game is over after $r$, so the win need not come after the reply.

\textbf{A.17 Remark.} Theorem~A.16 is a statement about $\mathrm{P}$ only. A won promotion in $\mathrm{P}$ need not be a
won chess game, so it does not give (V3) for chess, as the main text says. With Remark~A.12(b) and $Q$ as in Theorem~A.11, certificates for every
Z-move at every position of $\mathrm{R}(Q)$ give (V3) in $\mathrm{P}$.

\textbf{Further material.} The working note \texttt{kings/theory/LOOPY.md} in the code copy has further results and the records of the computer checks; its numbering differs from the numbering used here.

\subsection{Two Further Proofs for Section~S2}\label{S-kingslong}
% Section S12, second subsection: full proofs of the two statements of Section S2 whose proofs are not in S12.1
% (the theorem on outcomes with spare moves, and the local test for (V2)). Items are numbered A.18, A.19 by hand.

This subsection proves two statements of Section~\ref{S-kings}: Theorem~\ref{thm:loopy} on outcomes with spare moves
(Theorem~A.18a), and the frozen-king test for (V2) (Proposition~A.19 and Corollary~A.19a). They are Theorem~3.4 and
Proposition~2.11 with Corollary~2.11a of the working note \texttt{kings/theory/LOOPY.md} in the code copy. We use the
notation of Parts A.I--A.III and cite items A.1--A.17 of the previous subsection. Every further definition is given
here, and every proof is a hand proof. The one computed statement (in Remark~A.19b) is marked as such and is used
nowhere.

\subsubsection*{A.V Games with decisive terminals and spare moves}

\textbf{D-games.} A \emph{D-game} $G$ is a finite set of positions with, for each position $x$, a finite set $L(x)$ of
White options and a finite set $R(x)$ of Black options. An option is a position or one of two \emph{terminals}:
$\oplus$, ``White has won'', and $\ominus$, ``Black has won''. A move to a terminal is a \emph{decisive event}; a move of
$\sigma$ to the terminal of $\bar\sigma$ is a \emph{suicide}. Play from a position $x$ with a side $s$ to move
alternates, the mover choosing one of its options. A move to a terminal ends play with that winner, a player with no
option loses, and infinite play is a draw. A side $\sigma$ \emph{wins} $(x,s)$ if it has a strategy that ends every play
with a win for $\sigma$ after finitely many moves. At most one side wins (play the two strategies against each other);
if neither does, $(x,s)$ is a \emph{draw}. $G$ is \emph{short} if the directed graph on its positions, with an edge from
$x$ to every position in $L(x)\cup R(x)$, is acyclic; a move from $x$ to $x$ is a cycle. Then a sequence of moves
between positions, made by either side in any order, visits distinct positions, so it has at most $h-1$ moves, where $h$
is the number of positions of $G$; and the positions can be listed so that every option of a position that is itself a
position comes before it.

\textbf{Sums.} The sum $A_1\boxplus\dots\boxplus A_k$ of D-games has the tuples of positions as positions. A move is a
move in one coordinate, and a move to a terminal in any coordinate ends the whole play with that terminal: every
decisive event ends the game. The sum is a D-game (an option of a tuple is a terminal, or the tuple with one coordinate
replaced by an option of it), and it is short if every $A_i$ is, since a move between tuples moves one coordinate
to an earlier position in its listing and leaves the others fixed.

\textbf{Spare moves.} The king group $Z$ of Section~\ref{S-kings}, in which a side that can wait passes and every other
move loses, is modelled by a D-game, also called $Z$, with a single position $z$ and $L(z)\subseteq\{z,\ominus\}$,
$R(z)\subseteq\{z,\oplus\}$. White \emph{can wait} if $z\in L(z)$, and Black if $z\in R(z)$. The move $z\to z$ is a
\emph{pass}: it changes only the side to move. The other options of $Z$ are suicides, the moves of $Z$ that lose. Without
suicides, $Z$ is the D-game form of $\mathrm{on}$ if only White can wait, of $\mathrm{off}$ if only Black can, and of
$\mathrm{dud}$ if both can. For a short D-game $G$ we study $G\boxplus Z$; we write its position $(x,z)$ as $x$. It is not
short when a side can wait.

\textbf{Forcing.} Let $G$ be short. To test forcing from $(x,s)$, let the opponent $\bar\sigma$ choose at each of its
turns any move of $G$ or a pass (whether or not $\bar\sigma$ can wait in $Z$), and let $\sigma$ choose only moves of $G$.
Then $\sigma$ \emph{forces a decisive event against a passing opponent} from $(x,s)$ if $\sigma$ has a strategy under
which every such play reaches the terminal of $\sigma$ after finitely many moves. Since the opponent always has a pass,
$\sigma$ cannot win such a play by leaving it without a move.

\textbf{Predicates.} Let $G$ be short. Going through the positions of $G$ in an order that lists every option that is a
position before the position it is an option of, define for each position $x$,
first $\lambda_{\mathrm L}(x)$ and then $\lambda_{\mathrm R}(x)$:
\begin{itemize}
\item[] $\lambda_{\mathrm L}(x)$: some White option of $x$ is $\oplus$, or is a position $y$ with $\lambda_{\mathrm R}(y)$;
\item[] $\lambda_{\mathrm R}(x)$: $\lambda_{\mathrm L}(x)$ holds, and every Black option of $x$ is $\oplus$ or a
position $y$ with $\lambda_{\mathrm L}(y)$.
\end{itemize}
Define $\varrho_{\mathrm R}(x)$ and $\varrho_{\mathrm L}(x)$ in the same way with the colours exchanged: some Black option
of $x$ is $\ominus$ or a position $y$ with $\varrho_{\mathrm L}(y)$; and $\varrho_{\mathrm R}(x)$ holds and every White
option of $x$ is $\ominus$ or a position $y$ with $\varrho_{\mathrm R}(y)$. The subscript names the side to move
($\mathrm L$ for White, $\mathrm R$ for Black). By Lemma~A.18, $\lambda_s(x)$ says that White forces a decisive event
against a passing opponent from $(x,s)$, and $\varrho_s(x)$ the same for Black; this is the condition for winning when the
opponent can wait (Theorem~A.18a), but not in general. By definition $\lambda_{\mathrm R}(x)$ implies
$\lambda_{\mathrm L}(x)$, and $\varrho_{\mathrm L}(x)$ implies $\varrho_{\mathrm R}(x)$.

\textbf{A.18 Lemma (forcing and blocking).} \emph{Let $G$ be a short D-game, $Z$ as above, $x$ a position of $G$ and $s$
a side to move. (a) If $\lambda_s(x)$, White has a strategy in $G\boxplus Z$ that uses only moves of $G$ and wins
against every play of Black, passes included, whether or not Black can wait; in particular White forces a decisive
event against a passing opponent from $(x,s)$. (b) Let $\lambda_s(x)$ fail, and grant Black a pass at every Black turn
(a pass changes only the side to move). Then Black has a strategy that uses only moves of $G$ and passes, under which no
play reaches $\oplus$, whatever White does (moves of $G$, passes, suicides); in particular White does not force a
decisive event against a passing opponent from $(x,s)$. If Black can wait in $Z$, this is a strategy in $G\boxplus Z$
itself. (c) Parts (a) and (b) also hold with the colours exchanged, $\varrho$ in place of $\lambda$; in (b) White is
granted the pass, and the strategy is one in $G\boxplus Z$ if White can wait.}

\begin{proof}
(a) White keeps the invariant that at the current position $y$, with $t$ to move, $\lambda_t(y)$ holds; this is true at
the start. With White to move, $\lambda_{\mathrm L}(y)$ gives a White option that is $\oplus$, and White wins, or a
position $y'$ with $\lambda_{\mathrm R}(y')$; White plays it. With Black to move, $\lambda_{\mathrm R}(y)$ holds. A Black
pass leads to $y$ with White to move, and $\lambda_{\mathrm L}(y)$ holds since $\lambda_{\mathrm R}$ implies
$\lambda_{\mathrm L}$. A Black move of $G$ leads to $\oplus$ or to a position $y'$ with $\lambda_{\mathrm L}(y')$, by the
second clause of $\lambda_{\mathrm R}(y)$; a Black suicide in $Z$ is $\oplus$. So the invariant holds, no play reaches
$\ominus$, and White always has a move. Every White move is a move of $G$, to $\oplus$ or to a position, and the moves of
$G$ between positions in a play are at most $h-1$, since $G$ is short. White moves at every second turn, so every play
has at most $2h+1$ moves. It ends at $\oplus$ or with Black unable to move; either way White wins. Against a passing
opponent Black always has a move, so the play ends at $\oplus$.

(b) Black keeps the invariant that $\lambda_t(y)$ fails at the current position $y$ with $t$ to move. With Black to move,
$\lambda_{\mathrm R}(y)$ fails. If $\lambda_{\mathrm L}(y)$ fails, Black uses its granted pass. Otherwise the second
clause of $\lambda_{\mathrm R}(y)$ fails, so some Black option of $y$ is $\ominus$, and Black plays it and wins, or is a position
$y'$ where $\lambda_{\mathrm L}(y')$ fails, and Black plays it. With White to move, $\lambda_{\mathrm L}(y)$ fails, so no
White option of $y$ is $\oplus$ and every White option that is a position $y'$ has $\lambda_{\mathrm R}(y')$ false. A
White pass leads to $y$ with Black to move, where $\lambda_{\mathrm R}(y)$ fails because $\lambda_{\mathrm L}(y)$ does; a
White suicide is $\ominus$. So the invariant holds, and neither White nor Black (under this strategy) ever moves to
$\oplus$. Black's prescribed replies are legal in the forcing test with White forcing; every White move of $G$ keeps the
invariant, and if White to move has no option in $G$ the test ends without reaching $\oplus$. So White does not force a
decisive event against a passing opponent. If Black can wait in $Z$, the granted pass is the pass of $Z$, and the strategy is one in
$G\boxplus Z$.

(c) is (a) and (b) with the colours exchanged.
\end{proof}

\textbf{A.18a Theorem (outcomes with spare moves; Theorem~\ref{thm:loopy}).} \emph{Let $G$ be a short D-game, for
instance a sum of short D-games, let $Z$ be as above, and consider $G\boxplus Z$ from a position $x$ with $s$ to move.
\begin{itemize}
\item[(a)] If both sides can wait: White wins if and only if $\lambda_s(x)$, if and only if White forces a decisive event
against a passing opponent; Black wins if and only if $\varrho_s(x)$, if and only if Black forces a decisive event
against a passing opponent; otherwise $(x,s)$ is a draw. A side that wins has a winning strategy that never passes, and
having the move never hurts: $\lambda_{\mathrm R}(x)$ implies $\lambda_{\mathrm L}(x)$, and $\varrho_{\mathrm L}(x)$
implies $\varrho_{\mathrm R}(x)$.
\item[(b)] If only White can wait: every play is finite and $(x,s)$ is not a draw. Black wins if and only if
$\varrho_s(x)$, if and only if Black forces a decisive event against a passing White; otherwise White wins.
\item[(c)] If only Black can wait: the mirror image of (b).
\end{itemize}}

\begin{proof}
(a) If $\lambda_s(x)$, White wins by Lemma~A.18(a), with a strategy that uses only moves of $G$. If $\lambda_s(x)$
fails, Black follows the strategy of Lemma~A.18(b), which is a strategy in $G\boxplus Z$ since Black can wait; under it no
play reaches $\oplus$, and Black always has a move (a pass or the chosen option). A play ends in a White win only at
$\oplus$ or with Black unable to move, so White does not win. The two parts of Lemma~A.18 also show that $\lambda_s(x)$
holds exactly when White forces a decisive event against a passing opponent. The statements for Black follow in the same
way from Lemma~A.18(c), whose blocking strategy for White is a strategy in $G\boxplus Z$ since White can wait. If neither
$\lambda_s(x)$ nor $\varrho_s(x)$ holds, neither side wins, which is a draw: each side, by its blocking strategy, avoids
defeat against every play of the other, and when both follow these strategies the play never ends. The implications are
the remark after the definition of the predicates.

(b) Black cannot pass, so every Black move is a move of $G$ between positions, or a move to a terminal (in $G$ or a
suicide), which ends play. The moves of $G$ between positions in a play, by either side, are at most $h-1$, since passes
leave the position of $G$ unchanged and $G$ is short. So Black makes at most $h$ moves, and since the turns alternate
every play has at most $2h+1$ moves. Each position has finitely many options, so the game from $(x,s)$ is a finite game
tree whose every leaf has a winner (a terminal, or a player with no option), and by Zermelo's theorem one side
wins: there are no draws. If $\varrho_s(x)$, Black wins by Lemma~A.18(c), part (a) with colours exchanged; its strategy uses only moves of $G$ and
beats White's passes. If $\varrho_s(x)$ fails, White follows the strategy of Lemma~A.18(c), part (b) with colours
exchanged, which uses moves of $G$ and passes and is a strategy in $G\boxplus Z$ since White can wait. Under it no play
reaches $\ominus$ and White always has a move; every play is finite, so it ends at $\oplus$ or with Black unable to move, and
White wins. As in (a), $\varrho_s(x)$ holds exactly when Black forces a decisive event against a passing White.

(c) is (b) with the colours exchanged.
\end{proof}

\textbf{A.18b Remark (the statement of Section~\ref{S-kings}).} Theorem~\ref{thm:loopy} is (a) and (b): ``the
components form a short game in which every decisive event ends the game'' is $G=A_1\boxplus\dots\boxplus A_k$ with short
$A_i$; ``let a side that can wait pass at any time'' and ``let every other move in $Z$ lose'' are the D-game $Z$: a side
that can wait may pass at any of its turns, and every other move of $Z$ is a suicide, an immediate win for the opponent.
The forcing and blocking strategies above never prescribe a suicide, and each of them covers an opponent who plays one.
``Against a passing opponent'' means against an opponent who chooses at each turn between a move of $G$ and a pass, not
one who must pass: if $G$ has one position $x$ with $L(x)=\{\oplus\}$ and $R(x)=\{\ominus\}$, and White is to move, Black
reaches $\ominus$ against a White that always passes, yet $\varrho_{\mathrm L}(x)$ fails and White wins by moving to
$\oplus$. In (b) Black can win only at a terminal, since White always has a pass, so ``Black wins'' means ``Black forces $\ominus$''; White may also win because Black runs out of moves. The theorem is about the
D-game. That a king stepping back and forth in chess acts as such a pass is a further hypothesis, not treated here: the
steps must be legal at every position reached, and must change nothing else (in chess a king move also ends an e.p.\
right, while the pass of $Z$ keeps the position of $G$, right included).

\subsubsection*{A.VI A local test for (V2) with frozen kings}

\textbf{Setting.} We use the frozen-king game $\mathrm{F}$ of Part~A.III: chess in which the first promotion ends the game
and wins for the promoter, and a legal king move ends the game at once with a loss for its maker. Checkmate wins; a side
with no legal move (king moves count) that is not in check is stalemated, a draw; infinite play is a draw. The board is
flat or a cylinder with $H\ge5$ ranks. Let $y$ be a legal position (the kings not adjacent, the side not to move not in
check) with pawns and the two kings only, no promotion having occurred, with side $s$ to move and with the
e.p.\ right, if any, created by the double step that led to $y$ (Part~A.I). Fix a side $\beta$, the beneficiary, and let $\omega=\bar\beta$. For a pawn $q$ of either colour let $\rho(q)$ be its rank counted
from $\beta$'s side: $\rho(q)$ is its rank if $\beta$ is White and $H+1$ minus its rank if $\beta$ is Black. So $\beta$'s
pawns move to larger $\rho$, $\omega$'s to smaller $\rho$, and ``ahead of a pawn of $\beta$'' means larger $\rho$. For a
pawn $p$ of $\beta$ let
\[
d=H-\rho(p)\ \text{ if }\rho(p)\ge3,\qquad d=H-3\ \text{ if }\rho(p)=2,
\]
the number of moves $p$ needs to promote when it moves alone, with a double step first if $\rho(p)=2$; and let $e=d$ if
$s=\omega$ and $e=d-1$ if $s=\beta$. If play reaches the $d$-th move of $p$ and $p$ has moved at every turn of $\beta$,
then $\omega$ has made exactly $e$ moves before that move.

\textbf{The game $G_\omega(y;p)$.} Delete $p$ from $y$. Only $\omega$ moves, one move after another. $\beta$'s other
pawns stand still and may be captured; both kings stand still, occupy their squares and are never captured. $\omega$'s
pawns step, double step and capture by the pawn rules of Part~A.I, with the kings as obstacles and with no test of check,
and may capture e.p.\ only as the first move, only if $s=\omega$, and only with the right of $y$ (by (H1) that right
allows no capture of $p$). Let $\mu_\omega(y;p)$ be the least number of moves of a sequence in $G_\omega(y;p)$
after which some pawn of $\omega$ stands on its last rank, or $\infty$ if there is none. This minimum is over all pawns of
$\omega$, passed or not, and over all sequences, captures of $\beta$'s pawns included.

\textbf{Hypotheses.} Let $f$ be the file of $p$; the \emph{strip} is the set of files at distance at most $1$ from $f$.
\begin{itemize}
\item[(H1)] \emph{Clear path.} $p$ passes the e.p.-aware test at $y$ (Part~A.I): no pawn of $\omega$ in the strip has
$\rho>\rho(p)$, and $p$ has not just double stepped with a pawn of $\omega$ beside it. Moreover no pawn and no king of
either colour stands on a square of $f$ ahead of $p$, the promotion square included.
\item[(H2)] \emph{No check.} No pawn of either colour stands on a file at distance $1$ from the file of $\beta$'s king.
\item[(H3)] \emph{No stalemate.} $\omega$'s king has a neighbouring square $t$ at king distance (the maximum of the
file distance and the absolute difference of the ranks) at least $2$ from $\beta$'s king such that no pawn of either colour stands on
the file of $t$ or on a file at distance $1$ from it.
\item[(H4)] \emph{No interference.} Every pawn of $\beta$ other than $p$ on a file at distance $1$ from $f$ has
$\rho\le\rho(p)+1$.
\item[(H5)] \emph{Timing.} $\mu_\omega(y;p)\ge e+1$.
\end{itemize}

\textbf{A.19 Proposition (frozen-king passer test).} \emph{Under (H1)--(H5), $\beta$ wins $\mathrm{F}$ from $y$ with $s$
to move by the strategy $\Pi$: move $p$ forward at every turn, by a double step if $\rho(p)=2$. Every move of $\Pi$ is
legal when it is due, and every play in which $\beta$ follows $\Pi$ ends with a win for $\beta$, at the latest with the
$d$-th move of $p$, which promotes.}

\begin{proof}
Write $\rho_p$ for the current value of $\rho(p)$ and $\rho_0$ for its value at $y$. We consider plays in which $\beta$
follows $\Pi$ and $\omega$ moves arbitrarily, as long as the game has not ended.

\emph{Step 0 (frame).} No promotion has occurred, so the only pieces are pawns and kings, and by Lemma~A.1 a pawn changes
file only by capturing a pawn on the file it enters; a file with no pawn never gets one. A king move ends $\mathrm{F}$, so
in a play that goes on neither king has moved, and both stand where they stood at $y$. Under $\Pi$, $p$ is the only pawn
of $\beta$ that moves, and it stays on $f$.

\emph{Step 1 (checks and stalemate).} By (H2) and Step~0, no pawn ever stands on a file at distance $1$ from the file
of $\beta$'s king, so no pawn of $\omega$ attacks it, and $\omega$'s king, which does not move, is not adjacent to it
since $y$ is legal. So $\beta$ is never in check, and a move of $p$ is illegal only if its squares are occupied. By (H3)
and Step~0, the file of $t$ and the files at distance $1$ from it never hold a pawn; in particular $f$ is not one of them.
So $t$ is never occupied or attacked by a pawn, it is neither the square of $\beta$'s king nor adjacent to it, and the
move of $\omega$'s king to $t$ is legal at every turn of $\omega$, even in check. So $\omega$ is never stalemated.

\emph{Step 2 (invariant).} We show, by induction on the number of moves played, that
\begin{itemize}
\item[] (J$_\beta$) with $\beta$ to move, every pawn of $\omega$ in the strip has $\rho\le\rho_p+1$;
\item[] (J$_\omega$) with $\omega$ to move, every pawn of $\omega$ in the strip has $\rho\le\rho_p$, and $p$ cannot be
captured e.p.;
\item[] (I) no pawn and no king stands on a square of $f$ ahead of $p$.
\end{itemize}
At $y$ these are (H1): with $s=\omega$ it is (J$_\omega$) as stated, and with $s=\beta$ the bound $\rho\le\rho_0$ gives
(J$_\beta$). The kings never move (Step~0), so (I) concerns pawns only after $y$.

\emph{$\beta$ moves.} By (I) the squares of the move of $p$ are empty, and by Step~1 the move is legal; $p$ goes to
$\rho'\in\{\rho_p+1,\rho_p+2\}$ and stays on $f$, so (I) still holds. By (J$_\beta$) every pawn of $\omega$ in the strip
has $\rho\le\rho_p+1\le\rho'$. After a single step there is no right. After a double step, $\rho'=\rho_p+2$, every pawn
of $\omega$ in the strip is below $\rho'$, and a pawn that could capture $p$ e.p.\ would stand beside $p$ at $\rho'$; there
is none. So (J$_\omega$) holds.

\emph{$\omega$ moves a pawn.} A step or a double step lowers $\rho$ and keeps the file, so (J$_\beta$) holds for the
moved pawn. It cannot enter a square of $f$ ahead of $p$: a pawn of $\omega$ on $f$ has $\rho\le\rho_p$ by (J$_\omega$)
and moves to smaller $\rho$. A capture lands on the square of a captured pawn $q$ of $\beta$, or, e.p., on the square that
$q$ skipped. First, $q\ne p$: a capture of $p$ would need a pawn of $\omega$ in the strip at $\rho_p+1$, or a usable
right against $p$, and (J$_\omega$) excludes both. Second, a pawn $q\ne p$ of $\beta$ has not moved since $y$. If $q$ is on
$f$, it stands behind the square of $p$ at $y$ by (H1), and $p$ only advances; so the landing square, the square of $q$
or, e.p., the square behind it, is behind $p$, and (I) is kept. If $q$ is on a file at distance $1$ from $f$, then
$\rho(q)\le\rho_0+1\le\rho_p+1$ by (H4), and the landing rank is $\rho(q)$, or $\rho(q)-1$ for an e.p.\ capture. Captures
landing outside the strip do not concern (J$_\beta$). Third, an e.p.\ capture by $\omega$ needs a right created by
$\beta$; under $\Pi$ only $p$ double steps, and $p$ cannot be captured e.p., so any such capture is made at the first turn
from $y$, with $s=\omega$, using the right of $y$ against a pawn $q\ne p$; it is covered above. So after the move every pawn of
$\omega$ in the strip has $\rho\le\rho_p+1$, which is (J$_\beta$), and (I) holds. (Pawns of $\omega$ never move to
larger $\rho$, but a capture of a pawn of $\beta$ on a file next to $f$ brings a pawn of $\omega$ into the strip at the
rank of the captured pawn, or one rank lower after an e.p.\ capture; (H4) is what bounds it.)

\emph{Step 3 ($p$ survives and is never blocked).} By (J$_\omega$), at no turn of $\omega$ can $p$ be captured, by a
normal or an e.p.\ capture; by (I) its path is empty at every turn of $\beta$. With Step~1, every move of $\Pi$ is legal.

\emph{Step 4 ($\omega$'s pawn moves are moves of $G_\omega$).} Take a prefix of a play that follows $\Pi$ which contains no king
move and no promotion, except that its last move may be a promotion by a pawn of $\omega$. After each move of the prefix, the pieces other than $p$ stand on the same
squares as after the same pawn moves of $\omega$, in the same order, in $G_\omega(y;p)$ (a promoting pawn counted on its arrival square): the kings and $\beta$'s other
pawns have not moved, the same pawns of $\beta$ have been captured in both, and $p$ was never captured (Step~3). We
compare the moves in order. A step or double step of $\omega$ in the play uses empty squares, which are also empty in
$G_\omega$, where $p$ is deleted. A capture takes a pawn $q\ne p$ of $\beta$, which stands on the same square in
$G_\omega$. Any e.p.\ capture by $\omega$ uses the right of $y$ at $\omega$'s first move (Step~2), and the same capture,
by the same pawn, is allowed in $G_\omega$. So the pawn moves of $\omega$, in order, form a sequence of
$G_\omega(y;p)$, and if $\omega$ promotes with its $k$-th move then $k\ge\mu_\omega(y;p)$.

\emph{Step 5 (count).} Let a play follow $\Pi$. Before the $d$-th move of $p$, $\omega$ has at most $e$ turns. If
$\omega$ makes a king move at one of them, $\beta$ wins at once. Otherwise its moves there are pawn moves, and by Step~4
and (H5) none of them promotes, since $e<\mu_\omega(y;p)$. By Step~1, $\omega$ always has a legal move, the king move to
$t$, so it is neither stalemated nor checkmated, and $\beta$ is never in check; by Step~3, $\beta$ always has the move of
$\Pi$. So a play that has not ended with a king move of $\omega$ reaches the $d$-th move of $p$, which is legal by Step~3
and promotes, and $\beta$ wins.
\end{proof}

\textbf{A.19a Corollary ((V2) with frozen kings).} \emph{Let $Q$ be a king group with its context as in Part~A.III, let
$Y\in\mathrm{R}(Q)$ have side $M$ to move, and let $\mu$ be a context move at $Y$ that produces a passed pawn and is
adjudicated for $\sigma$ in $W$. Let $y'$ be the position after $\mu$, with $N=\bar M$ to move. If for some pawn $p$ of
$\sigma$ the hypotheses (H1)--(H5) hold at $y'$ with $s=N$ and $\beta=\sigma$ (so $e=d$ if $\sigma=M$ and $e=d-1$ if
$\sigma=N$), then $\sigma$ wins $y'$ in $\mathrm{F}$. If this holds for every such move at every $Y\in\mathrm{R}(Q)$, then
(V2) holds with chess replaced by $\mathrm{F}$, the form used in Remark~A.12(b).}

\begin{proof}
$y'$ is a legal position reached by a legal move. No promotion has occurred: play from $Q$ to $Y$ contains no decisive
event, hence no Z-move and no context promotion (proof of Theorem~A.11), and $\mu$ does not promote, by Lemma~A.2(c) as
in that proof. So Proposition~A.19 applies at $y'$ with $N$ to move and $\beta=\sigma$.
\end{proof}

\textbf{A.19b Remark (what the test gives and what it does not).} (a) The test is a sufficient condition. For any
fixed race rule, the corollary applies the test to the side that $W$ credits at that move; a different rule can change
the side to be tested. (H5) counts every pawn of $\omega$, passed or not, and every sequence of $\omega$'s moves with $p$
removed, including captures of $\beta$'s other pawns and any e.p.\ capture available at the first move; a count of
passed pawns alone does not suffice. (b) It is a statement about $\mathrm{F}$, where king moves lose; Step~0 uses that the kings never move. It gives
neither $\mathrm{P}$ nor chess. On $\cyl{8}$ with six ranks, with White king a1 and pawn c2, Black king a6 and White to
move, (H1)--(H5) hold for $\beta$ White and $p=$\,c2 ($\mu_\omega=\infty$), so White wins $\mathrm{F}$; yet in
$\mathrm{P}$, with real kings, the position is a draw \lab{C1}. Thus, when the test passes at
every move that Corollary~A.19a quantifies over, the corollary supplies (V2) for Theorem~A.11 in the form of Remark~A.12(b) with $\mathrm{F}$, and not
for $\mathrm{P}$ or chess.
